%% =====================================================================
%%  ROUND 5, REVISED IN ROUNDS 6-9 -- the proof fragment: phases,
%%  overlap periodicity, and the mult-1 wall.  Paper-voice draft of
%%  what is PROVED, what is machine-verified-but-unproved, and exactly
%%  what remains for "rev is not L-reachable".
%%
%%  ROUND 6 REVISIONS:
%%   - Fact (no consecutive left-moves): REFUTED at larger sizes
%%     (B^inf-structured triple + the staircase).  Retained as the
%%     record of the small-domain verdict, marked refuted.
%%   - The staircase remark added: LDS ~ |w|/2 at mult 3 for a fixed
%%     two-pass expression -- Conjecture A (LDS <= C(E) mult) REFUTED.
%%   - Base case restated as a CONJECTURE (the round-5 derivation used
%%     the refuted fact; the statement has no counterexample).
%%
%%  ROUND 9 REVISIONS:
%%   - Conjecture basecase and the mult-1 theorem (Theorem B): REFUTED
%%     by the period-3 family E = [eps/tail^4 X].[tail(X)/ab]X on
%%     w = (bba)^k: LDS = k-1 at mult EXACTLY 1, k = 5..16, every row
%%     content-cross-checked.  The supply pigeonhole of
%%     Remark rem:textlevel does NOT generalize (2-char sigma, period
%%     3: the residue class is exactly the right size).
%%   - Proposition phasebound SURVIVES as the structure theorem of
%%     mult-1 outputs; its constant is value-dependent.
%%   - The atom route is DEAD as a separating program; what remains
%%     is the influence-crossing side and the descending-bijection
%%     invariant (new).
%%
%%  Machine verification (re-runnable):
%%    docs/proof/research/scratch/rev/verify_round5_phases.py
%%  (Lemma phasing: 4,000 copy-gap texts; Lemma overlap: 128,520 word
%%   pairs exhaustive; base case: 205,585 real pipeline instances.)
%%    docs/proof/research/scratch/rev/verify_round6_witness.py
%%  (wall refutations: the B^inf triple standalone + the staircase
%%   family to j=40, every row content-cross-checked; Conjecture A
%%   refutation.)
%%    docs/proof/research/scratch/rev/verify_round6_mult1_sweep.py
%%  (mult-1 hunt: no row with mult=1 and LDS>=4.)
%%    docs/proof/research/scratch/rev/verify_round7_mult1.py
%%  (Proposition phasebound: 1,157,184 exhaustive + 20,000 random,
%%   zero violations.)
%%    docs/proof/research/scratch/rev/verify_round9_supply.py
%%  (extended hunt: 581,856 pipelines, max LDS at mult 1 = 7 -- the
%%   refutation; death breakdown; residue check.)
%%    docs/proof/research/scratch/rev/verify_round9_witness.py
%%  (the witness family k = 5..16, per-row den cross-check; the
%%   descending-bijection probe: none with |w| >= 5.)
%% =====================================================================

\subsection{The Structure of a Shaving Pass}

Fix a value $A$ (the replacement of an earlier pass) occurring as
$m$ disjoint copies in the current text, and a deletion pass
$[\,\epsilon/B\,]$.  A copy is entered at an \emph{offset} $o$: either
$0$ (the scan stands at the copy's first atom) or the landing point of
a match that straddled the left boundary; in the latter case $B$'s
suffix of length $o$ is $A$'s prefix of length $o$.  A copy is exited
at a \emph{depth} $j$: either no match crosses the right boundary
($j = 0$) or a single match does, consuming $j$ copy atoms, and then
$B$'s prefix of length $j$ is $A$'s suffix of length $j$.  Write
\begin{align*}
    O(A,B) &\triangleq \{\, o \in [1, \min(|A|,|B|{-}1)] :
              A[{:}o] = B[|B|{-}o{:}] \,\}, \\
    J(A,B) &\triangleq \{\, j \in [1, \min(|A|,|B|{-}1)] :
              A[|A|{-}j{:}] = B[{:}j] \,\}.
\end{align*}

\begin{lemma}
    (\emph{Phases})
    \label{lem:phase}
    The residual of a copy is a function of $(o, j)$ alone; the possible
    pairs lie in $(\{0\} \cup O(A,B)) \times (\{0\} \cup J(A,B))$, plus one
    further class of residuals, the empty one, produced by matches that
    cover a copy entirely.  Hence the copies' residuals take at most
    $(1 + |O(A,B)|)(1 + |J(A,B)|) + 1$ distinct values.
\end{lemma}

\begin{proof}
    From offset $o$ the greedy scan through the copy is deterministic
    and independent of the surrounding text while its windows lie
    inside $A$; the only interaction with the following gap is the
    single match that crosses the right boundary, at depth $j$, after
    which the scan has left the copy.  A match ending at atom $o$
    covers $A[0{:}o]$ followed by the preceding text, so $B$ ends with
    $A[{:}o]$; a match crossing from atom position $q$ covers
    $A[q{:}]$ followed by gap characters, so $B$ starts with
    $A[q{:}]$ and $j = |A| - q$.
\end{proof}

This supersedes the cruder bound
$|B| \cdot |\Sigma|^{|B|-1}$: the per-pass diversity of shaving is
governed not by the needle's length but by its \emph{overlap sets} with
the copied value.  Those sets are classical objects:

\begin{lemma}
    (\emph{Overlap Periodicity})
    \label{lem:overlap}
    Let $o_1 < o_2 < \cdots < o_k$ be the elements of $O(A,B)$ and
    $o_{\max} = o_k$.  Then:
    (i) every $o_i$, $i < k$, is a border of the prefix $A[{:}o_{\max}]$
    (its length-$o_i$ prefix equals its length-$o_i$ suffix);
    (ii) every difference $o_k - o_i$ is a period of $A[{:}o_k]$;
    (iii) if $O(A,B)$ is an arithmetic progression with step $d$, then
    $d$ is a period of $A[{:}o_{\max}]$: the straddled prefix is an
    $(o_{\max}/d)$-fold repetition of a word of length $d$.  The mirror
    statements hold for $J(A,B)$ on $A$'s suffix.
\end{lemma}

\begin{proof}
    All elements of $O$ are suffixes of $B$ of the corresponding
    lengths and prefixes of $A$ of the same lengths; nested suffixes of
    one word give the border and period relations by the standard
    arguments (a word whose first $b$ characters equal its last $b$ has
    period $|w| - b$).  \emph{Verified exhaustively} for all $A$ with
    $|A| \leq 8$ and $B$ with $|B| \leq 7$ over $|\Sigma| = 2$
    ($128{,}520$ pairs), all clauses.
\end{proof}

\noindent The quantitative form one might hope for -- $|O| \leq
o_{\max}/p$ for the minimal period $p$ -- is \emph{false} in general
($A = \mathtt{aabaabaa}$, $B = \mathtt{xaabaa}$ gives $O = \{1,2,5\}$
at $p = 3$); the border-chain structure of (i)--(ii) is the correct
statement.  In the extremal case -- the run mechanism of the
variable-depth constructions -- $O$ is a full arithmetic progression
and the copied value's prefix is a repetition, which is exactly why
uniform runs are where all known disorder lives, and why that disorder
is content-invisible.

\subsection{The Left-Move Wall}

Call a pair of copies $c < c'$ (in text order) with residuals $R_c
\ni p$, $R_{c'} \ni q$ a \emph{left-move} if $q < p$: the later copy
keeps a \emph{smaller} position.  At multiplicity $1$ the residuals
are pairwise disjoint on atoms, and every strictly decreasing chain of
w-atoms through the copies induces left-moves whenever the copied
value's provenance is increasing.  The machine verdict:

%% ------------------------------------------------------------------
%% ROUND 6 REVISION.  The Fact below, as stated in round 5, is REFUTED
%% at larger sizes.  It is retained (struck through) as the record of
%% what the small domain suggested, followed by the correct mechanism.
%% ------------------------------------------------------------------

\begin{fact}
    (\emph{No Consecutive Left-Moves}; REFUTED in round 6)
    \label{fact:leftmove}
    \emph{Round 5 statement}: after one deletion pass, among copies
    whose residuals are pairwise disjoint, the longest strictly
    decreasing position chain has length exactly $2$ -- verified
    exhaustively over all texts with $|A| \leq 4$, $|B| \leq 3$, gaps
    of length $\leq 2$, at most three copies ($1{,}152{,}480$ texts;
    $4{,}730$ under the all-disjoint convention, $23{,}368$ under the
    exists-a-disjoint-pair convention), and on $40{,}000$ random texts.
    \emph{Round 6 refutation}: a targeted search over
    $B^{\infty}$-structured texts (the habitat pinned by the power
    equations) yields the triple $B = \mathtt{ababa}$, $A =
    \mathtt{bababab}$, gaps $\mathtt{aa},\mathtt{a},\mathtt{a},
    \mathtt{a}$ with residuals $\{4\}, \{2\}, \{0,6\}$ and picks
    $4 > 2 > 0$; and the staircase family of
    Remark~\ref{rem:staircase} realizes chains of length
    $\sim |w|/2$ through pairwise disjoint singleton residuals.  The
    round-5 domain was too small.
\end{fact}

The correct mechanism is the \emph{modular staircase}.  When the
needle $B$ is a factor of the same periodic stream that the copies
tile, the greedy scan matches at spacing $m + \pi/2$ (where $\pi$ is
the stream's period) and emits a fixed small number of atoms between
consecutive matches; emission $c$ lands at offset $(m +
\pi/2)\,c \bmod (n + g)$ in its copy unit, which \emph{descends by a
constant step per emission}.  The strictly descending run before the
first wrap has length linear in $|w|$; the wrap and the end-of-text
dump each duplicate labels.  This is what long decreasing chains look
like, and why they come with multiplicity $\geq 2$.

\begin{remark}
    (\emph{The Staircase}; machine-verified)
    \label{rem:staircase}
    For $E = [\,\epsilon/\mathrm{init}^2 X\,]\,[\,X/b\,]X$ and $w_j =
    \mathtt{b}(\mathtt{ab})^{j}$, the output provenance is
    \[
        (n-2,\; n-4,\; \ldots,\; 3,\; 1,\; n-2,\; n-2,\; n-1),
        \qquad n = 2j+1
    \]
    (the descent is always on \emph{odd} labels; $n$ is odd),
    with $\mathrm{LDS} = j$ \emph{exactly} and multiplicity $3$
    \emph{exactly}, for $j = 2, \ldots, 40$ (and $\mathrm{LDS} =
    j/2 + 1$ at multiplicity $4$--$5$ for $\mathrm{init}^4$).  Hence
    $\mathrm{LDS}/\mathrm{mult} \to \infty$ for a \emph{fixed}
    two-pass expression: the bound
    $\mathrm{LDS} \leq C(E) \cdot \mathrm{mult}$ is false.  Every row
    is cross-checked against an independent implementation of the
    denotation.
\end{remark}

\begin{remark}
    (\emph{The Period-3 Family}; machine-verified -- the mult-$1$
    wall does not exist)
    \label{rem:period3}
    For $E = [\,\epsilon/\mathrm{tail}^4 X\,]\,[\,\mathrm{tail}(X)/
    \mathtt{ab}\,]X$ and $w_k = (\mathtt{bba})^{k}$, the output
    provenance is
    \[
        (0,\; 1,\; n-3,\; n-6,\; n-9,\; \ldots,\; 3,\; n-2,\; n-1),
        \qquad n = 3k,
    \]
    with $\mathrm{LDS} = k - 1$ and multiplicity $1$ \emph{exactly},
    for $k = 5, \ldots, 16$ (every row cross-checked against the
    independent denotation).  The mechanism is the modular staircase
    of the preceding paragraph with the period $\pi = 3$ and a
    two-character $\sigma = \mathtt{ab}$: the emission offsets descend
    by $s = 3$, so the staircase is \emph{confined to one residue
    class} mod $3$ -- the $\mathtt{b}$-positions $0, 3, \ldots, 3k-3$
    of $w$, of which $w$ supplies exactly $k-1$: one per copy,
    exactly enough, never revisited.  The interior gap atoms (the
    same residue class) are consumed by the matches; the text ends on
    the trailing gap, so no tail is dumped; the leading and trailing
    gap labels lie outside the staircase's values.  With this the
    mult-$1$ theorem is \emph{false}: $\mathrm{LDS}$ is unbounded at
    multiplicity $1$, even for a fixed two-pass expression.  The
    supply pigeonhole of Remark~\ref{rem:textlevel} does not
    generalize: for period $2$ and one-character $\sigma$ the
    relevant residue class of $w$ is too small (that is the
    pigeonhole); for period $3$ and two-character $\sigma$ it is
    exactly the right size.
\end{remark}

%% ------------------------------------------------------------------
%% ROUND 9 REVISION.  The Conjecture below -- and with it the mult-1
%% theorem ("Theorem B": mult 1 => LDS <= C(E)) -- is REFUTED by the
%% period-3 family of Remark rem:period3.  Retained as the record of
%% what rounds 5-8's evidence suggested, followed by the refutation.
%% ------------------------------------------------------------------

\begin{conjecture}
    (\emph{Base Case}; restated round 6 -- REFUTED in round 9)
    \label{prop:basecase}
    \emph{Round 6 statement}: for the two-pass fragment
    $[\,A/\sigma\,]\,[\,\epsilon/B\,]$ with $A =
    \llbracket R \rrbracket(w)$: at multiplicity $1$,
    \[
        \mathrm{LDS}(\mathrm{prov}) \;\leq\; 2 \cdot
        \mathrm{LDS}(\mathrm{prov}_A) + 1 .
    \]
    \emph{Round 9 refutation} (Remark~\ref{rem:period3}): $E =
    [\,\epsilon/\mathrm{tail}^4 X\,]\,[\,\mathrm{tail}(X)/
    \mathtt{ab}\,]X$ on $w = (\mathtt{bba})^{k}$ has
    $\mathrm{LDS}(\mathrm{prov}_A) = 1$ ($A = \mathrm{tail}(w)$
    increasing) and $\mathrm{LDS} = k - 1$ at multiplicity $1$
    \emph{exactly}, for $k = 5, \ldots, 16$ (every row
    content-cross-checked).  The rounds 5--8 evidence -- $205{,}585$
    instances, then $355{,}928$, then $581{,}856$ realizable
    pipelines -- was a domain artifact: the refuting family enters
    the earlier sweeps only at $k \geq 5$ ($|w| \geq 15$), where its
    $\mathrm{LDS} = k - 1 \geq 4$ first stands out.
\end{conjecture}

\begin{proposition}
    (\emph{Phase Bound at Multiplicity $1$}; proved)
    \label{prop:phasebound}
    For the two-pass fragment $[\,A/\sigma\,]\,[\,\epsilon/B\,]$: if the
    output provenance has multiplicity $1$, then
    \[
        \mathrm{LDS}(\mathrm{prov}) \;\leq\;
        \Bigl(\bigl(1 + |O(A,B)|\bigr)\bigl(1 + |J(A,B)|\bigr) +
        1\Bigr) \cdot \mathrm{LDS}(\mathrm{prov}_A) \;+\; 1 .
    \]
\end{proposition}

\begin{proof}
    Multiplicity $1$ permits at most one \emph{surviving copy per
    realized phase}: two surviving copies with the same entry/exit
    pair $(o,j)$ have identical residuals (Lemma~\ref{lem:phase}); a
    nonempty residual then shares an offset, hence a label
    $\mathrm{prov}_A[i]$, with the other copy, forcing multiplicity
    $\geq 2$.  So the number of surviving copies is at most the number
    of realized nonempty residual classes, which
    Lemma~\ref{lem:phase} bounds by $(1+|O|)(1+|J|)+1$.  Each
    surviving copy contributes at most $\mathrm{LDS}(\mathrm{prov}_A)$
    to a decreasing chain: within one copy, text order is ascending
    offsets, so the picks of a chain restricted to that copy are a
    decreasing subsequence of $\mathrm{prov}_A$.  The gap survivors
    are the $\sigma$-interiors of $w$; their labels strictly increase
    in text order, so at most one of them lies on any decreasing chain.
\end{proof}

\noindent \emph{Verified} (round 7): $1{,}157{,}184$ exhaustive
pipeline instances ($|A| \leq 4$ with label sequences of LDS $1,2,3$,
$|B| \leq 3$, $8$ sigmas, $|w| \leq 6$; $591{,}315$ at multiplicity
$1$) and $20{,}000$ random larger ones: zero violations of both
claims, via an independently written tagged simulator whose output
content is cross-checked against the reference denotation.

The proposition's constant is \emph{value-dependent}: $|O|, |J|$ grow
with the values $A, B$.  The remaining gap to a uniform mult-$1$
theorem is to show that multiplicity $1$ bounds the number of
\emph{realized} phases by $C(E)$ alone -- and here the round-7
finding is that any proof \emph{must} use realizability:

\begin{remark}
    (\emph{The text-level statement is false; realizability is
    load-bearing})
    \label{rem:textlevel}
    As a statement about texts with free injective labels, the bound
    $2\,\mathrm{LDS}(\mathrm{prov}_A)+1$ is \emph{false}.  With $A =
    \mathtt{babababab}$, $B = \mathtt{abababa}$, gaps
    $\mathtt{a},\mathtt{a},\mathtt{a},\mathtt{a},\mathtt{a}$, the
    residuals are $\{6\},\{4\},\{2\},\{0,8\}$ and the chain $6 > 4 >
    2 > 0$ has length $4$; over aligned alternating texts, chains up
    to length $|A|/2$ (observed: $6$) are realized at multiplicity
    $1$.  None of these lifts to a pipeline: an alternating $A$ of
    odd length $n$ needs $(n{+}1)/2$ input positions of one parity,
    but the $\sigma$-site structure of any $w$ realizing such a text
    (with $\sigma = \mathtt{b}$, one site per such position) supplies
    only $(n{-}1)/2$; a nondecreasing $\mathrm{prov}_A$ cannot
    bridge the gap (adjacent repeats clash parity, non-adjacent
    repeats force constant stretches).  This supply pigeonhole is a
    sketch, and it does \emph{not} generalize: the period-$3$ family
    of Remark~\ref{rem:period3} is the counterexample (the residue
    class is exactly the right size there).
\end{remark}

\begin{lemma}
    (\emph{Border--Period}; classical, machine-verified)
    \label{lem:borderperiod}
    If $x$ has $t \geq 2$ borders and minimal period $p$, then
    $p\,(t-1) \;\leq\; |x| - 1$.
\end{lemma}

\noindent \emph{Verified} exhaustively over all binary strings $|x|
\leq 11$.  With Lemma~\ref{lem:phase} this quantifies the periodicity
half of the mult-$1$ theorem: if $K$ phases are realized, then
$\max(|O|, |J|) \geq \sqrt{K} - 1$, so $A$ has $\geq \sqrt{K}-1$
borders at one end and hence a period $\leq (|A|-1)/(\sqrt{K}-2)$
there.  Few-phases-or-periodic is now a theorem.  The supply half,
however, is where the program ends: the period-$3$ family of
Remark~\ref{rem:period3} is periodic at \emph{both} ends, embeds
into $w$ along a single residue class of exactly the right size, and
realizes $k-1$ phases at multiplicity $1$ -- the periodicity and
supply constraints do not combine into a uniform bound.  The direct
sweeps confirm the escalation in domain: $355{,}928$ realizable
pipelines (round 8) never exceed $\mathrm{LDS} = 2$ at multiplicity
$1$; $581{,}856$ over the period-$3$ families with phase variants and
3-character $\sigma$ (round 9) reach $\mathrm{LDS} = 7$ -- the
witness of Remark~\ref{rem:period3}, at $|w| = 24$.

\subsection{What This Gives, and What Remains}

What survives rounds 5--9 is a \emph{structure theorem}, not a
complexity bound: at multiplicity $1$, every decreasing chain through
the output takes at most one atom per realized phase, one per copy,
and at most one gap atom (Proposition~\ref{prop:phasebound}) -- the
witness family of Remark~\ref{rem:period3} satisfies it, realizing
$k-1$ phases and $k-1$ chain links.  What does \emph{not} survive is
the separation program built on it: Conjecture~A
($\mathrm{LDS} \leq C(E)\cdot\mathrm{mult}$) fell in round 6
(Remark~\ref{rem:staircase}); Conjecture~B (mult $1 \Rightarrow$
$\mathrm{LDS} \leq C(E)$) and the base case
(Conjecture~\ref{prop:basecase}) fell in round 9
(Remark~\ref{rem:period3}).  There is no $(\mathrm{mult},
\mathrm{LDS})$-separating invariant: $\mathcal{L}$ realizes
multiplicity $1$ with $\mathrm{LDS}$ linear in $|w|$, so reversal's
profile (mult $1$, $\mathrm{LDS} = n$) is \emph{not} rare in
$\mathcal{L}$.  The atom route, as a program of bounding
$\mathrm{LDS}$ by $\mathrm{mult}$ and expression constants, is dead.

Two live directions remain, and both were visible before round 9 --
the refutation sharpens their standing:

\begin{enumerate}
    \item \emph{The influence-crossing route} (the quantitative form
    of open problem~2 in Remark~\ref{rem:price}): reversal is
    length-preserving with $\chi = \binom{n}{2}$ crossings, while
    every length-preserving pipeline examined so far has
    $\chi = O(|w|)$.  The length-discard rule -- the caveat that any
    crossing bound must apply to length-changing witnesses, where
    $\chi = 0$ vacuously -- remains load-bearing, and the
    period-$3$ family is itself length-changing ($|w| = 3k$,
    $|\mathrm{out}| = k+3$), so this route is untouched by the
    refutation.  A hypothetical constant-gate witness (output rebuilt
    from constants selected by content) is likewise invisible here;
    crossings are the only side that sees both.
    \item \emph{The descending-bijection invariant} (new in round 9):
    reversal's provenance is a \emph{bijection} (every input position
    survives exactly once: mult $1$, $|\mathrm{prov}| = |w|$, all
    labels present) \emph{and} fully descending ($\mathrm{LDS} =
    |\mathrm{prov}|$).  The two-pass fragment provably cannot be a
    descending bijection: a copy with $\geq 2$ surviving atoms has
    ascending labels, so a full descent keeps at most one atom per
    copy, while a bijection forbids deletion -- the constraints are
    incompatible.  Machine probe over the round-7 exhaustive domain
    ($1{,}157{,}184$ instances at the \emph{text} level, a strict
    superset of realizable pipelines since $\mathrm{prov}_A$ is an
    arbitrary injective labeling): descending bijections exist only
    for $|w| \leq 4$ ($110$ essential families, all requiring label
    patterns that constants do not carry); none with $|w| \geq 5$.
    The period-$3$ family is far from a bijection
    ($|\mathrm{prov}| = k+3$ vs $|w| = 3k$: deletion-heavy).  The
    question is whether deeper pipelines can close the gap.
\end{enumerate}

The three sibling obstructions stay orthogonal, as
Remark~\ref{rem:price} records: the once-hinge and direction-hinge
canonical functions pass all of the above with $\chi = 0$.
